\chapter[Metody Fouriera i rozdzielenie zmiennych]{Szeregi Fouriera, rozdzielenie zmiennych i zagadnienia Sturma--Liouville'a}
\label{chap:pde-fourier}

W rozdziałach~\ref{chap:pde-heat} i~\ref{chap:pde-wave} uzyskaliśmy rozwiązania na
całej przestrzeni i w prostych obszarach. Na przedziale lub prostokącie
bardziej naturalne jest korzystanie z \emph{modów własnych} operatora
przestrzennego. Metoda rozdzielenia zmiennych będzie skuteczna tylko wtedy,
gdy uzasadnimy, że funkcje własne tworzą dostatecznie bogatą rodzinę.
Najpierw przeprowadzimy odpowiedni rachunek dla szeregów Fouriera,
a następnie zbadamy jego wersję dla regularnych zagadnień
Sturma--Liouville'a. Formalny rachunek dystrybucji przedstawimy w rozdziale~\ref{chap:pde-distributions},
a transformatę Fouriera na $\R^n$ i dystrybucje temperowane ---
w rozdziale~\ref{chap:pde-fourier-Rd}, po wprowadzeniu
przestrzeni Schwartza. Rozróżnienie szeregów Fouriera na torusie
i transformaty na całej przestrzeni jest istotne także dla normalizacji.

\section{Szereg Fouriera: współczynniki, ortogonalność i zbieżność}
\label{sec:pde6-fourier}

Przez $\mathbb T=\R/(2\pi\mathbb Z)$ rozumiemy okrąg z
współrzędną $x\in[-\pi,\pi]$ i utożsamionymi końcami.
Funkcję z $L^1(-\pi,\pi)$ utożsamiamy z jej
$2\pi$-okresowym przedłużeniem, określonym z dokładnością
do zbioru miary zero; wartości na końcach przedziału
nie wpływają na całki. Dla takiej funkcji definiujemy
$\widehat f(k)=(2\pi)^{-1}\int_{-\pi}^{\pi}f(x)e^{-ikx}\,\mathrm d\lambda_1(x)$,
$k\in\mathbb Z$, oraz sumę częściową
$S_Nf(x)=\sum_{|k|\le N}\widehat f(k)e^{ikx}$.
W tym rozdziale $L^2(\mathbb T)$ oznacza zespoloną przestrzeń Hilberta
z iloczynem skalarnym
$\langle f,g\rangle=(2\pi)^{-1}\int_{-\pi}^{\pi}f(x)\overline{g(x)}\,\mathrm d\lambda_1(x)$.
Pojęcie całki Lebesgue'a i podstawowe własności $L^p$ traktujemy
jako znane; szerszą teorię funkcjonalną rozwiniemy w dalszych tomach.
Zwracamy uwagę na normalizację: w tym rozdziale norma $L^2(\T)$
pochodzi od miary $\mathrm d\lambda_1(x)/(2\pi)$.
W dalszych rozdziałach o PDE na torusie stosujemy zwykle
nieznormalizowaną miarę $\mathrm d\lambda_d(x)$; wzory
na wartości własne i współczynniki Fouriera pozostają zgodne,
ale norma względem miary nieznormalizowanej w wymiarze $d$
jest $(2\pi)^{d/2}$ razy większa od normy względem
miary unormowanej do masy $1$ (na torusie
o okresie $2\pi$ w każdej współrzędnej).
Na odcinku $(0,L)$ używamy natomiast zwykłej
miary Lebesgue'a. Ponieważ miara unormowana na
$\mathbb T$ ma masę $1$, nierówność Cauchy'ego--Schwarza
daje $(2\pi)^{-1}\|f\|_{L^1(-\pi,\pi)}\leq\|f\|_2$,
więc współczynniki Fouriera są określone dla każdego
$f\in L^2(\mathbb T)$.

\begin{lemat}[Ortogonalność i nierówność Bessela]
\label{lem:pde6-bessel}
Niech $f\in L^2(\mathbb T)$ oraz $N\geq0$ będzie całkowite. Wówczas
$\langle e^{ikx},e^{i\ell x}\rangle=\delta_{k\ell}$,
$\|S_Nf\|_2^2=\sum_{|k|\le N}|\widehat f(k)|^2$ i
$\sum_{|k|\le N}|\widehat f(k)|^2\le\|f\|_2^2$.
\end{lemat}
\begin{proof}
Dla $k=\ell$ całka z jedynki względem miary
unormowanej wynosi $1$. Jeśli $k\ne\ell$, to
\[
 \frac1{2\pi}\int_{-\pi}^{\pi}e^{i(k-\ell)x}\,
 \mathrm d\lambda_1(x)
 =\frac{e^{i(k-\ell)\pi}-e^{-i(k-\ell)\pi}}
 {2\pi i(k-\ell)}=0,
\]
ponieważ $k-\ell$ jest całkowite. Otrzymujemy
ortonormalność. Rozwijając kwadrat normy skończonej
sumy $S_Nf$ i korzystając z tej ortonormalności,
dostajemy $\|S_Nf\|_2^2=\sum_{|k|\le N}|\widehat f(k)|^2$.
Ponadto dla $|\ell|\leq N$
\[
 \langle f-S_Nf,e^{i\ell x}\rangle
 =\widehat f(\ell)-\sum_{|k|\leq N}
 \widehat f(k)\delta_{k\ell}=0.
\]
Zatem $S_N$ jest rzutem ortogonalnym na przestrzeń
rozpiętą przez te wykładniki. W szczególności
$\langle f-S_Nf,S_Nf\rangle=0$, skąd
\[
 \|f\|_2^2=\|S_Nf\|_2^2+\|f-S_Nf\|_2^2.
\]
Nieujemność ostatniego składnika daje nierówność Bessela.
\end{proof}

Nie można jednak wywnioskować z ortogonalności, że $S_Nf\to f$.
Do tego potrzebujemy zupełności układu wykładniczego.

\begin{lemat}[Jądro Fejéra]
\label{lem:pde6-fejer}
Dla całkowitego $N\geq0$ niech
\begin{equation}\label{eq:pde6-fejer}
 F_N(t)=\frac1{N+1}\left|\sum_{j=0}^Ne^{ijt}\right|^2
 =\frac1{N+1}\left(\frac{\sin((N+1)t/2)}{\sin(t/2)}\right)^2
\end{equation}
(w punktach $t\in2\pi\mathbb Z$ przyjmujemy
$F_N(t)=N+1$, zgodnie z pierwszą postacią wzoru). Wówczas $F_N\ge0$,
$(2\pi)^{-1}\int_{-\pi}^{\pi}F_N(t)\,\mathrm d\lambda_1(t)=1$ oraz
dla każdego $\delta\in(0,\pi)$
$\int_{\delta\le|t|\le\pi}F_N(t)\,\mathrm d\lambda_1(t)\to0$.
\end{lemat}
\begin{proof}
Nieujemność jest bezpośrednia z definicji.
Dla $t\notin2\pi\mathbb Z$ z sumy geometrycznej
$\sum_{j=0}^N e^{ijt}=(1-e^{i(N+1)t})/(1-e^{it})$
i tożsamości $|1-e^{i\alpha}|=2|\sin(\alpha/2)|$
otrzymujemy drugą postać~\eqref{eq:pde6-fejer}.
Rozwijając kwadrat modułu, mamy
\[
 F_N(t)=\frac1{N+1}\sum_{j,\ell=0}^N e^{i(j-\ell)t}
 =\sum_{|k|\leq N}\left(1-\frac{|k|}{N+1}\right)e^{ikt}.
\]
Dla ustalonego $|k|\leq N$ istnieje dokładnie
$N+1-|k|$ par $(j,\ell)$ z różnicą $j-\ell=k$:
dla $k\geq0$ wybieramy $\ell=0,\ldots,N-k$,
a dla $k<0$ zamieniamy role indeksów.
Wyraz stały jest równy $1$, a pozostałe
wykładniki mają całkę zero na pełnym okresie,
skąd normalizacja. Gdy $\delta\leq|t|\leq\pi$,
mamy $|\sin(t/2)|\geq\sin(\delta/2)>0$.
Dlatego
\[
 \int_{\delta\leq|t|\leq\pi}F_N(t)\,\mathrm d\lambda_1(t)
 \leq\frac{2\pi}{(N+1)\sin^2(\delta/2)}
 \longrightarrow0.
\]
\end{proof}

\begin{twierdzenie}[Fejér i zupełność układu trygonometrycznego]
\label{thm:pde6-fejer}
Niech $f\in C(\mathbb T)$. Średnie Cesàra
$\sigma_N f=(N+1)^{-1}\sum_{j=0}^N S_jf$
zbieżne są jednostajnie do $f$ na $\mathbb T$.
W konsekwencji wielomiany trygonometryczne są gęste
w $C(\mathbb T)$ w normie supremum i w $L^2(\mathbb T)$
w normie $L^2$.
\end{twierdzenie}
\begin{proof}
W średniej z sum częściowych współczynnik
przy $\widehat f(k)e^{ikx}$ wynosi
$(N+1-|k|)/(N+1)$ dla $|k|\leq N$, ponieważ
wyraz ten występuje dla $j=|k|,\ldots,N$.
Z drugiej strony podstawienie $y=x-t$ daje
\[
 \frac1{2\pi}\int_{-\pi}^{\pi}e^{ikt}f(x-t)\,
 \mathrm d\lambda_1(t)=e^{ikx}\widehat f(k).
\]
Po podstawieniu całkujemy po odcinku
$[x-\pi,x+\pi]$, który możemy zastąpić
$[-\pi,\pi]$, gdyż $f(y)e^{-iky}$ jest
$2\pi$-okresowa i całka po pełnym okresie
nie zależy od jego początku (przez podział
odcinka i przesunięcie o wielokrotność okresu).
Rozwinięcie skończone z lematu~\ref{lem:pde6-fejer}
daje więc
\[
 \sigma_Nf(x)=\frac1{2\pi}\int_{-\pi}^{\pi}
 F_N(t)f(x-t)\,\mathrm d\lambda_1(t).
\]
Odejmujemy $f(x)$ i dzielimy całkę na $|t|<\delta$
oraz $\delta\le|t|\le\pi$. Z jednostajnej ciągłości $f$
możemy dobrać $\delta$ tak, aby pierwszy składnik
miał moduł co najwyżej $\varepsilon$ dla wszystkich $x,N$.
Drugi ma moduł nie większy niż
$2\|f\|_\infty(2\pi)^{-1}\int_{|t|\ge\delta}F_N(t)
\,\mathrm d\lambda_1(t)$, co dąży do zera.
Ponieważ $F_N$ ma unormowaną całkę $1$, pierwszy
składnik jest rzeczywiście ograniczony przez
$\varepsilon$, niezależnie od $N$. Zatem dla
dostatecznie dużych $N$ norma supremum różnicy
nie przekracza $2\varepsilon$. Dowodzi to
jednostajnej zbieżności.

Uzasadnijmy również gęstość w $L^2$ z uwzględnieniem
okresowości. Funkcje proste są gęste w $L^2$
na przedziale skończonej miary. Dla mierzalnego
$A\subset(-\pi,\pi)$ oraz $\eta>0$ regularność
miary Lebesgue'a daje zbiór zwarty $K\subset A$
i otwarty $O\supset A$, $O\subset(-\pi,\pi)$,
takie że $\lambda_1(O\setminus K)<\eta$.
Gdy $K$ jest niepusty, wybieramy
$0<r<\operatorname{dist}(K,\R\setminus O)$.
Funkcja
\[
 \varphi(x)=\max\{0,1-\operatorname{dist}(x,K)/r\}
\]
jest ciągła, równa $1$ na $K$, ma zwarty nośnik
zawarty w $O$ i wartości w $[0,1]$.
Jeśli $K=\varnothing$, bierzemy $\varphi=0$.
W obu przypadkach
$\|\chi_A-\varphi\|_2^2\leq\eta/(2\pi)$.
Przedłużenie $\varphi$ okresowo jest ciągłe,
bo $\varphi$ znika w otoczeniu końców przedziału.
Aproksymując w ten sposób każdy składnik
funkcji prostej i stosując nierówność trójkąta,
otrzymujemy gęstość $C(\mathbb T)$ w $L^2(\mathbb T)$.
Dla $f\in L^2$ i $\varepsilon>0$ wybieramy
$g\in C(\mathbb T)$ z $\|f-g\|_2<\varepsilon/2$,
a następnie wielomian trygonometryczny $p$
z $\|g-p\|_\infty<\varepsilon/2$.
Ponieważ $\|g-p\|_2\leq\|g-p\|_\infty$,
mamy $\|f-p\|_2<\varepsilon$.
\end{proof}

\begin{twierdzenie}[Zbieżność w $L^2$ i tożsamość Parsevala]
\label{thm:pde6-parseval}
Dla dowolnego $f\in L^2(\mathbb T)$ zachodzą
$S_Nf\to f$ w normie $L^2$ oraz
\begin{equation}\label{eq:pde6-parseval}
 \frac1{2\pi}\int_{-\pi}^{\pi}|f(x)|^2\,\mathrm d\lambda_1(x)
 =\sum_{k\in\mathbb Z}|\widehat f(k)|^2.
\end{equation}
\end{twierdzenie}
\begin{proof}
Dla $\varepsilon>0$ wybieramy wielomian trygonometryczny $p$
z $\|f-p\|_2<\varepsilon$. Dla $N$ nie mniejszych niż stopień
$p$ mamy $S_Np=p$. Ponieważ $f-S_Nf$ jest
prostopadła do przestrzeni, do której należy
$S_Nf-p$, zachodzi
\[
 \|f-p\|_2^2=\|f-S_Nf\|_2^2+\|S_Nf-p\|_2^2.
\]
Zatem $\|f-S_Nf\|_2\leq\|f-p\|_2<\varepsilon$,
co dowodzi zbieżności.
W konsekwencji $\|S_Nf\|_2\to\|f\|_2$,
a z tożsamości $\|S_Nf\|_2^2=
\sum_{|k|\leq N}|\widehat f(k)|^2$
w lemacie~\ref{lem:pde6-bessel} otrzymujemy
\eqref{eq:pde6-parseval}, ponieważ sumy
nieujemnych wyrazów rosną do pełnej sumy po $\mathbb Z$.
\end{proof}

\begin{uwaga}[Jakiego rodzaju zbieżność została udowodniona?]
Sama zbieżność w normie $L^2$ nie pociąga
za sobą zbieżności jednostajnej ani punktowej
całego ciągu. Nie rozstrzygamy tutaj odrębnego
zagadnienia zbieżności prawie wszędzie sum
częściowych Fouriera. Dla funkcji ciągłych jednostajnie zbieżne są
\emph{średnie Fejéra}, niekoniecznie zwykłe sumy $S_N$.
W dalszej części do różniczkowania szeregów użyjemy
gładkich danych, dla których zachodzi dodatkowe szybkie
zanikanie współczynników.
\end{uwaga}

\begin{lemat}[Zanikanie współczynników gładkiej funkcji]
\label{lem:pde6-decay}
Jeśli $f\in C^m(\mathbb T)$, $m\geq1$ całkowite, to dla $k\ne0$
$|\widehat f(k)|\le\|f^{(m)}\|_{L^1(-\pi,\pi)}/(2\pi|k|^m)$.
Jeśli $f\in C^\infty(\mathbb T)$, to szereg Fouriera wraz
ze wszystkimi szeregami otrzymanymi przez różniczkowanie
zbiega jednostajnie.
\end{lemat}
\begin{proof}
Dla $k\ne0$ całkowanie przez części daje
\[
 \widehat f(k)=\frac1{ik}\widehat{f'}(k),
\]
ponieważ $f(\pi)e^{-ik\pi}=f(-\pi)e^{ik\pi}$.
Powtarzając ten rachunek $m$ razy (pochodne
do rzędu $m-1$ mają zgodne wartości na końcach
przedziału), dostajemy
$\widehat f(k)=(ik)^{-m}\widehat{f^{(m)}}(k)$,
co po oszacowaniu modułu całki daje żądaną nierówność.

Dla $f\in C^\infty$ i ustalonego całkowitego
$r\geq0$ wybieramy $m\geq r+2$. Wówczas
\[
 \sum_{k\ne0}|k|^r|\widehat f(k)|
 \leq\frac{\|f^{(m)}\|_{L^1(-\pi,\pi)}}{2\pi}
 \sum_{k\ne0}|k|^{r-m}<\infty.
\]
Kryterium Weierstrassa daje bezwzględną
i jednostajną zbieżność szeregu
$\sum_{k\in\mathbb Z}(ik)^r\widehat f(k)e^{ikx}$,
przy czym dla $r=0$ wyraz $k=0$ oznacza
$\widehat f(0)$, a dla $r>0$ wynosi zero.
Oznaczmy jego ciągłą sumę przez $G_r$.
Z całkowania przez części także dla $k=0$
(jeśli $r>0$, całka z $f^{(r)}$ po okresie
jest zerowa) współczynniki tego szeregu
są dokładnie współczynnikami Fouriera $f^{(r)}$.
Twierdzenie~\ref{thm:pde6-parseval} daje więc
zbieżność tych samych sum częściowych do $f^{(r)}$
w normie $L^2$. Z jednostajnej zbieżności
mamy zarazem zbieżność w $L^2$ do $G_r$.
Jednoznaczność granicy w $L^2$ daje
$G_r=f^{(r)}$ prawie wszędzie, a ponieważ
obie funkcje są ciągłe, równość zachodzi
w każdym punkcie: niezerowa różnica
w jednym punkcie byłaby niezerowa na
pewnym otwartym łuku dodatniej miary.
W szczególności $G_0=f$ i wszystkie
różniczkowania wyraz po wyrazie prowadzą
do właściwych pochodnych.
\end{proof}

\section{Zagadnienie własne laplasjanu na odcinku}
\label{sec:pde6-dirichlet}

Dla $L>0$ rozważmy
\begin{equation}\label{eq:pde6-eigen}
 -X''(x)=\lambda X(x),\quad 0<x<L,\qquad X(0)=X(L)=0.
\end{equation}
\begin{twierdzenie}[Pełne spektrum Dirichleta w jednym wymiarze]
\label{thm:pde6-spectrum}
Dla $L>0$ zagadnienie~\eqref{eq:pde6-eigen} ma
nietrywialne rozwiązanie $X\in C^2([0,L])$
wtedy i tylko wtedy, gdy
$\lambda=\lambda_n=(n\pi/L)^2$ dla pewnego całkowitego $n\ge1$.
Jego przestrzeń rozwiązań jest rozpięta przez
$X_n(x)=\sin(n\pi x/L)$. Ponadto funkcje
$e_n(x)=\sqrt{2/L}\sin(n\pi x/L)$, $n\ge1$,
tworzą bazę ortonormalną w $L^2(0,L)$.
\end{twierdzenie}
\begin{proof}
Mnożąc~\eqref{eq:pde6-eigen} przez $\overline X$
i całkując przez części otrzymujemy
$\lambda\int_0^L|X(x)|^2\,\mathrm d\lambda_1(x)
=\int_0^L|X'(x)|^2\,\mathrm d\lambda_1(x)$.
Ponieważ $X$ nie jest identycznie zerowa, całka z $|X|^2$
jest dodatnia. Gdyby całka z $|X'|^2$ była zerowa,
ciągłość $X'$ wymuszałaby $X'=0$ na całym odcinku.
Wtedy $X$ byłaby stała, a z $X(0)=0$ wynikałoby $X=0$,
sprzecznie z założeniem. Stąd $\lambda$ jest rzeczywista
i dodatnia. Ogólne rozwiązanie ODE to
$A\sin(\sqrt\lambda x)+B\cos(\sqrt\lambda x)$.
Warunek w $0$ daje $B=0$, a nietrywialność wymusza $A\ne0$.
Warunek w $L$ jest więc równoważny
$\sqrt\lambda L=n\pi$ dla pewnego całkowitego $n\ge1$.
Odwrotnie, bezpośrednie różniczkowanie $X_n$ daje
$-X_n''=\lambda_n X_n$ i oba warunki brzegowe.

Dla $m,n\ge1$ tożsamość
$2\sin a\sin b=\cos(a-b)-\cos(a+b)$ daje
\[
 \int_0^L\sin(m\pi x/L)\sin(n\pi x/L)\,
 \mathrm d\lambda_1(x)=\frac L2\delta_{mn}.
\]
Istotnie, całka z $\cos(j\pi x/L)$ znika dla każdego
niezerowego całkowitego $j$, a dla $j=0$ wynosi $L$.
Stąd rodzina $(e_n)$ jest ortonormalna.

Dla zupełności niech $f\in L^2(0,L)$. Przedłużamy ją
nieparzyście do $(-L,L)$, otrzymując $\widetilde f$.
Mamy $\int_{-L}^L|\widetilde f|^2\,\mathrm d\lambda_1
=2\int_0^L|f|^2\,\mathrm d\lambda_1$.
Podstawienie $x=L\theta/\pi$ pozwala zastosować
twierdzenie~\ref{thm:pde6-parseval} do funkcji
$\theta\mapsto\widetilde f(L\theta/\pi)$.
Współczynniki wykładnicze na okresie $2L$ są równe
\[
 c_k=\frac1{2L}\int_{-L}^L\widetilde f(x)
 e^{-ik\pi x/L}\,\mathrm d\lambda_1(x).
\]
Nieparzystość daje $c_0=0$, $c_{-n}=-c_n$ oraz
$c_n=-\frac{i}{L}\int_0^Lf(x)\sin(n\pi x/L)\,
\mathrm d\lambda_1(x)$ dla $n\ge1$.
Nie wymaga to rzeczywistości funkcji $f$.
Grupując wyrazy o indeksach $n,-n$, otrzymujemy
\[
 \sum_{|k|\le N}c_ke^{ik\pi x/L}
 =\sum_{n=1}^N\langle f,e_n\rangle_{L^2(0,L)}e_n(x).
\]
Lewa strona zbiega do $\widetilde f$ w $L^2(-L,L)$;
ograniczenie do $(0,L)$ nie zwiększa normy błędu.
Zatem prawa strona zbiega do $f$ w $L^2(0,L)$,
co dowodzi zupełności. Z ortonormalności i przejścia
do granicy norm sum częściowych wynika również
$\|f\|_{L^2(0,L)}^2=\sum_{n\ge1}|\langle f,e_n\rangle|^2$.
\end{proof}

\begin{wniosek}[Nierówność Poincarégo dla funkcji gładkich]
\label{cor:pde6-poincare}
Jeżeli $v\in C^1([0,L])$, $v(0)=v(L)=0$, to
\begin{equation}\label{eq:pde6-poincare}
 \int_0^L|v(x)|^2\,\mathrm d\lambda_1(x)
 \le \left(\frac L\pi\right)^2
 \int_0^L|v'(x)|^2\,\mathrm d\lambda_1(x).
\end{equation}
Stała $(L/\pi)^2$ jest optymalna.
\end{wniosek}
\begin{proof}
Niech $v_n=\langle v,e_n\rangle_{L^2(0,L)}$.
Oznaczmy $d_n(x)=\sqrt{2/L}\cos(n\pi x/L)$
dla $n\ge1$. Całkowanie przez części daje
\[
 \langle v',d_n\rangle
 = [v d_n]_0^L-\int_0^L v(x)d_n'(x)\,
 \mathrm d\lambda_1(x)
 = (n\pi/L)v_n,
\]
gdyż $v$ znika na obu końcach i $d_n'=-(n\pi/L)e_n$.
Z tożsamości $2\cos a\cos b=\cos(a-b)+\cos(a+b)$
i całek kosinusów obliczonych w poprzednim dowodzie
wynika ortonormalność $(d_n)_{n\ge1}$.
Nierówność Bessela dla tej rodziny (nie potrzebujemy
jej zupełności) i tożsamość Parsevala dla bazy sinusowej dają
\[
 \|v'\|_2^2\ge\sum_{n\ge1}(n\pi/L)^2|v_n|^2
 \ge(\pi/L)^2\sum_{n\ge1}|v_n|^2
 =(\pi/L)^2\|v\|_2^2.
\]
To jest~\eqref{eq:pde6-poincare}. Dla
$v(x)=\sin(\pi x/L)$ mamy $\|v\|_2^2=L/2$ i
$\|v'\|_2^2=(\pi/L)^2 L/2$, więc żadna mniejsza
stała nie może być poprawna dla wszystkich $v$.
\end{proof}

\begin{uwaga}[Po co zachowywać stałe w oszacowaniach?]
Nierówność~\eqref{eq:pde6-poincare} wzmacnia
oszacowanie elementarne uzyskane w rozdziale~\ref{chap:pde-heat}.
Najmniejsza wartość własna $\lambda_1=\pi^2/L^2$
wyznacza dokładną szybkość zaniku pierwszego modu
równania ciepła i najniższą częstość drgań struny.
\end{uwaga}

\section{Rozdzielenie zmiennych: równanie ciepła}
\label{sec:pde6-heat}

Niech $\kappa,L>0$, a $u_t=\kappa u_{xx}$ dla $t>0$,
$x\in(0,L)$, z $u(t,0)=u(t,L)=0$ i $u(0,x)=g(x)$.
Podstawienie $u(t,x)=T(t)X(x)$ daje
$T'X=\kappa TX''$. Dla nietrywialnego iloczynu wybieramy
$t_*$ z $T(t_*)\ne0$ i kładziemy
$\lambda=-T'(t_*)/(\kappa T(t_*))$. Otrzymujemy
$-X''=\lambda X$ dla wszystkich $x$, wraz z zerowymi
warunkami Dirichleta. Wybór $x_*$ z $X(x_*)\ne0$
daje następnie $T'=-\kappa\lambda T$ dla wszystkich $t$.
Nie dzielimy przez funkcje w ich punktach zerowych.
Dla $\lambda=\lambda_n$ otrzymujemy
$T(t)=T(0)e^{-\kappa\lambda_nt}$; stałą $T(0)$ włączamy
do amplitudy modu.


\begin{lemat}[Jądro Gaussa i kontrakcja na odcinku]
\label{lem:pde6-dirichlet-heat}
Niech $L,\kappa>0$ i niech
$g\in C([0,L])$ spełnia $g(0)=g(L)=0$.
Oznaczmy przez $\widetilde g$ jej nieparzyste
$2L$-okresowe przedłużenie na $\R$. Dla $t>0$ połóżmy
\[
 (Q_tg)(x)=\int_{\R}(4\pi\kappa t)^{-1/2}
 e^{-z^2/(4\kappa t)}\widetilde g(x-z)
 \,\mathrm d\lambda_1(z),\qquad x\in[0,L].
\]
Wówczas $Q_tg$ znika dla $x=0,L$,
$\|Q_tg\|_{C([0,L])}\le\|g\|_{C([0,L])}$ oraz
$\|Q_tg-g\|_{C([0,L])}\to0$ dla $t\downarrow0$.
Ponadto współczynnik sinusowy $Q_tg$ odpowiadający
$\sin(n\pi x/L)$ jest równy
$e^{-\kappa(n\pi/L)^2t}b_n$, gdzie
$b_n=(2/L)\int_0^Lg(x)\sin(n\pi x/L)\,\mathrm d\lambda_1(x)$.
\end{lemat}
\begin{proof}
Przedłużenie jest ciągłe i jednostajnie ciągłe na $\R$,
a jego norma supremum jest równa normie $g$.
Oznaczmy jądro Gaussa w definicji przez $G_t(z)$.
Jest ono nieujemne, a podstawienie $z=2\sqrt{\kappa t}\,r$
i całka Gaussa dają $\int_{\R}G_t\,\mathrm d\lambda_1=1$.
Stąd bezwzględna zbieżność całki i kontrakcja.
Ten sam wzór określa funkcję na całej prostej.
Jednostajna ciągłość $\widetilde g$ i normalizacja jądra
dają ciągłość splotu. Okresowość $\widetilde g$ daje
$Q_tg(x+2L)=Q_tg(x)$, a podstawienie $z\mapsto-z$
i parzystość $G_t$ dają $Q_tg(-x)=-Q_tg(x)$.
W szczególności $Q_tg(0)=0$ i
$Q_tg(L)=Q_tg(-L)=-Q_tg(L)$.

Dla $\varepsilon>0$ dobieramy $\delta>0$ z jednostajnej
ciągłości tak, aby dla $|z|<\delta$ różnica
$|\widetilde g(x-z)-\widetilde g(x)|$ była mniejsza
niż $\varepsilon$ dla wszystkich $x$. Wówczas
\[
 \|Q_tg-g\|_\infty\le\varepsilon+2\|g\|_\infty
 \int_{|z|\ge\delta}G_t(z)\,\mathrm d\lambda_1(z).
\]
Podstawienie $z=2\sqrt{\kappa t}\,r$ zamienia ostatnią
całkę na $\pi^{-1/2}\int_{|r|\ge\delta/(2\sqrt{\kappa t})}
e^{-r^2}\,\mathrm d\lambda_1(r)$, która dąży do zera
z całkowalności Gaussa. Ponieważ $\varepsilon$ jest
dowolne, otrzymujemy zbieżność jednostajną.

Niech $K_t(z)=\sum_{j\in\mathbb Z}G_t(z+2jL)$.
Dla $z\in[-L,L]$ i $|j|\ge1$ mamy
$|z+2jL|\ge(2|j|-1)L$, więc szereg zbiega jednostajnie
na okresie dla ustalonego $t>0$ przez kryterium Weierstrassa.
Podział prostej na przedziały długości $2L$ i okresowość
$\widetilde g$ dają
\[
 Q_tg(x)=\int_{-L}^L K_t(z)\widetilde g(x-z)\,
 \mathrm d\lambda_1(z).
\]
Dla współczynników $c_n(f)=(2L)^{-1}
\int_{-L}^L f(x)e^{-in\pi x/L}\,\mathrm d\lambda_1(x)$
całkowanie szeregu i podstawienie na każdym przedziale dają
\[
 2L c_n(K_t)=\int_{\R}G_t(z)e^{-in\pi z/L}\,
 \mathrm d\lambda_1(z).
\]
Czynnik $2L$ wynika z tego, że splot zapisano względem
zwykłej miary Lebesgue'a. Z twierdzenia Fubiniego
(ograniczoność $\widetilde g$ i całkowalność jądra)
oraz podstawienia $y=x-z$ w całce po pełnym okresie
wynika $c_n(Q_tg)=2L c_n(K_t)c_n(\widetilde g)$.

Dla $\xi\in\R$ niech
$M_t(\xi)=\int_{\R}G_t(z)e^{-i\xi z}\,\mathrm d\lambda_1(z)$.
Różniczkowanie pod znakiem całki jest uzasadnione
całkowalnością $|z|G_t(z)$. Ponieważ
$zG_t(z)=-2\kappa tG_t'(z)$, całkowanie przez części
z zanikającym wyrazem brzegowym daje
\[
 M_t'(\xi)=2i\kappa t\int_{\R}G_t'(z)e^{-i\xi z}
 \,\mathrm d\lambda_1(z)=-2\kappa t\xi M_t(\xi).
\]
Warunek $M_t(0)=1$ daje $M_t(\xi)=e^{-\kappa t\xi^2}$.
Stąd $c_n(Q_tg)=e^{-\kappa(n\pi/L)^2t}c_n(\widetilde g)$.
Przejście od współczynników wykładniczych funkcji
nieparzystych do sinusowych, jak w dowodzie
twierdzenia~\ref{thm:pde6-spectrum}, daje tezę również
dla zespolonej funkcji $g$.
\end{proof}

\begin{twierdzenie}[Rozwiązanie przez szereg sinusowy]
\label{thm:pde6-heat}
Niech $\kappa,L>0$ i $g\in C([0,L])$ spełnia
$g(0)=g(L)=0$. Definiujemy
$b_n=(2/L)\int_0^Lg(x)\sin(n\pi x/L)\,\mathrm d\lambda_1(x)$.
Wówczas
\begin{equation}\label{eq:pde6-heat}
 u(t,x)=\sum_{n=1}^\infty b_n e^{-\kappa\lambda_nt}
 \sin(n\pi x/L),\qquad t>0,\;x\in[0,L],
\end{equation}
jest klasycznym rozwiązaniem równania ciepła dla $t>0$,
ciągłym na $[0,\infty)\times[0,L]$ i przyjmuje dane $g$
w sposób jednostajny, gdy $t\downarrow0$. Rozwiązanie to
jest jedyne w klasie funkcji ciągłych na każdym
$[0,T]\times[0,L]$, klasy $C^{1,2}$ we wnętrzu
i spełniających dane i warunki Dirichleta.
\end{twierdzenie}
\begin{proof}
Z ograniczoności $g$ mamy $|b_n|\le2\|g\|_\infty$.
Dla $t\ge\delta>0$ szereg i wszystkie szeregi
z pochodnymi czasowymi i przestrzennymi zbieżne są
jednostajnie, bo dla dowolnych $r,s\ge0$
$\sum_{n\ge1}n^{r+2s}e^{-\kappa(n\pi/L)^2\delta}<\infty$.
Równanie i warunki brzegowe sprawdzamy wyraz po wyrazie.
Dla $t>0$ całkowanie jednostajnie zbieżnego szeregu
i ortogonalność sinusów dają współczynniki sinusowe
$b_ne^{-\kappa\lambda_nt}$. Z lematu~\ref{lem:pde6-dirichlet-heat}
te same współczynniki ma $Q_tg$. Zupełność układu
sinusowego (twierdzenie~\ref{thm:pde6-spectrum})
daje równość w $L^2(0,L)$, a ciągłość obu funkcji
daje równość w każdym punkcie $[0,L]$.
Zatem $u(t,\cdot)=Q_tg$ i
$\|u(t,\cdot)-g\|_\infty\to0$ dla $t\downarrow0$.
Definiujemy $u(0,x)=g(x)$. Ciągłość w $(0,x_0)$
wynika z oszacowania
$|u(t,x)-g(x_0)|\le\|u(t,\cdot)-g\|_\infty+|g(x)-g(x_0)|$.
Dla dodatnich czasów wynika ze zbieżności jednostajnej
na pasach $t\ge\delta$. Nie zakładaliśmy punktowej
zbieżności zwykłych sum częściowych szeregu Fouriera $g$.

Jednoznaczność wynika z parabolicznej zasady maksimum
z rozdziału~\ref{chap:pde-heat}, zastosowanej do różnicy
i jej przeciwieństwa na każdym skończonym przedziale
czasowym. Dla funkcji zespolonych stosujemy ten argument
osobno do części rzeczywistej i urojonej.
\end{proof}

\begin{wniosek}[Ostry wykładnik zaniku w normie $L^2$]
\label{cor:pde6-decay}
Jeżeli $g\in C([0,L])$ i $g(0)=g(L)=0$, to dla
rozwiązania~\eqref{eq:pde6-heat}
\begin{equation}\label{eq:pde6-decay}
 \|u(t,\cdot)\|_{L^2(0,L)}
 \le e^{-\kappa\pi^2t/L^2}\|g\|_{L^2(0,L)},\qquad t\ge0.
\end{equation}
Stała w wykładniku jest optymalna.
\end{wniosek}
\begin{proof}
Dla $t>0$ stosujemy ortogonalność sinusów i Parsevala:
\[
\|u(t)\|_2^2=\frac L2\sum_{n\ge1}|b_n|^2e^{-2\kappa\lambda_nt}
\le e^{-2\kappa\lambda_1t}\|g\|_2^2.
\]
Dla $t=0$ jest to równość. Przy $g(x)=\sin(\pi x/L)$
zachodzi równość dla każdego $t$.
\end{proof}


\begin{stwierdzenie}[Rozwiązanie cieplne dla danych $L^2$]
\label{prop:pde6-L2-semigroup}
Niech $\kappa,L>0$, $g\in L^2(0,L)$ i niech
$g_n=\langle g,e_n\rangle_{L^2(0,L)}$ dla ortonormalnych
funkcji $e_n$ z twierdzenia~\ref{thm:pde6-spectrum}.
Szereg $u(t)=\sum_{n\ge1}e^{-\kappa\lambda_nt}g_ne_n$
wyznacza funkcję $u\in C([0,\infty);L^2(0,L))$,
przy czym $u(0)=g$, a dla każdego $t>0$ otrzymujemy
funkcję gładką względem $x\in[0,L]$, klasyczną
w $t>0$ i spełniającą zerowe warunki Dirichleta.
Zachodzi także $\|u(t)\|_2\le e^{-\kappa\lambda_1t}\|g\|_2$.
\end{stwierdzenie}
\begin{proof}
Z twierdzenia Parsevala mamy $\sum_{n\ge1}|g_n|^2<\infty$.
Dla $t\ge0$ szereg zbiega w $L^2$, a
\[
 \|u(t)-g\|_2^2=\sum_{n\ge1}
 |e^{-\kappa\lambda_nt}-1|^2|g_n|^2\longrightarrow0
 \quad(t\downarrow0)
\]
na mocy twierdzenia o zbieżności zmajoryzowanej
dla szeregów, ponieważ każdy składnik jest ograniczony
przez $|g_n|^2$. Dla dowolnej chwili $t_0\ge0$
stosujemy ten sam argument do
$\sum_{n\ge1}|e^{-\kappa\lambda_nt}-e^{-\kappa\lambda_nt_0}|^2
|g_n|^2$, z tą samą sumowalną majorantą.
Otrzymujemy ciągłość $L^2$ w czasie.
Dla $t\ge\delta>0$ nierówność Cauchy'ego--Schwarza
zapewnia jednostajną zbieżność szeregów pochodnych:
\[
 \sum_{n\ge1}|g_n|n^k e^{-\kappa\lambda_n\delta}
 \le\Bigl(\sum_{n\ge1}|g_n|^2\Bigr)^{1/2}
 \Bigl(\sum_{n\ge1}n^{2k}e^{-2\kappa\lambda_n\delta}\Bigr)^{1/2}
 <\infty.
\]
Dla pochodnej $r$ razy po $x$ i $s$ razy po $t$
stosujemy to oszacowanie z $k=r+2s$ oraz stałym
czynnikiem zależnym od $L,\kappa,r,s$. Kryterium
Weierstrassa daje jednostajną zbieżność wszystkich
szeregów pochodnych, w tym szeregu bez pochodnych.
Jego suma punktowa reprezentuje granicę w $L^2$,
gdyż zbieżność jednostajna na odcinku implikuje
zbieżność w $L^2$. Różniczkowanie wyraz po wyrazie
daje równanie i warunki brzegowe dla $t>0$.
Nie przypisujemy danym $L^2$ wartości punktowych
w chwili $t=0$. Wreszcie z Parsevala wynika
$\|u(t)\|_2^2=\sum_{n\ge1}e^{-2\kappa\lambda_nt}|g_n|^2
\le e^{-2\kappa\lambda_1t}\|g\|_2^2$.
\end{proof}

\section{Rozdzielenie zmiennych: równanie falowe i energia modów}
\label{sec:pde6-wave}

Rozważmy $u_{tt}=c^2u_{xx}$ na $(0,L)$,
z zerowymi warunkami Dirichleta oraz danymi
$u(0,x)=g(x)$, $u_t(0,x)=h(x)$.
Podstawiając $u(t,x)=T_n(t)X_n(x)$, gdzie
$X_n(x)=\sin(n\pi x/L)$ i $\lambda_n=(n\pi/L)^2$,
otrzymujemy $(T_n''+c^2\lambda_nT_n)X_n=0$.
Ponieważ $X_n$ nie jest identycznie zerowa, mamy
$T_n''+c^2\lambda_nT_n=0$ dla wszystkich $t$.
Oznaczmy $\omega_n=c\sqrt{\lambda_n}>0$. Rozwiązanie
tego równania z danymi $T_n(0)=a_n$, $T_n'(0)=b_n$
ma postać $a_n\cos(\omega_nt)+(b_n/\omega_n)\sin(\omega_nt)$,
co sprawdzamy przez różniczkowanie i podstawienie $t=0$.

\begin{twierdzenie}[Struna: reprezentacja modalna]
\label{thm:pde6-wave}
Niech $c,L>0$ i $g,h\in C_c^\infty(0,L)$.
Niech $a_n=(2/L)\int_0^Lg(x)\sin(n\pi x/L)
\,\mathrm d\lambda_1(x)$ i
$b_n=(2/L)\int_0^Lh(x)\sin(n\pi x/L)\,\mathrm d\lambda_1(x)$.
Wówczas funkcja
\begin{equation}\label{eq:pde6-wave}
 u(t,x)=\sum_{n=1}^\infty
 \left[a_n\cos(c\sqrt{\lambda_n}t)
 +\frac{b_n}{c\sqrt{\lambda_n}}
 \sin(c\sqrt{\lambda_n}t)\right]\sin(n\pi x/L)
\end{equation}
jest jedynym rozwiązaniem klasy $C^2(\R\times[0,L])$
powyższego zagadnienia wśród rozwiązań klasycznych.
Jego energia falowa
$E(t)=\frac12\int_0^L(|u_t(t,x)|^2+c^2|u_x(t,x)|^2)
\,\mathrm d\lambda_1(x)$ jest stała.
\end{twierdzenie}
\begin{proof}
Nieparzyste $2L$-okresowe przedłużenia $g,h$
są klasy $C^\infty$, ponieważ obie funkcje
zerują się w otoczeniach końców przedziału.
Po przeskalowaniu okresu lemat~\ref{lem:pde6-decay}
i zależność między współczynnikami wykładniczymi
i sinusowymi z dowodu twierdzenia~\ref{thm:pde6-spectrum}
dają dla każdego całkowitego $m\ge0$ oszacowanie
$|a_n|+|b_n|\le C_m n^{-m}$ z pewną stałą $C_m$.
Niech $A_n(t)=a_n\cos(\omega_nt)
+(b_n/\omega_n)\sin(\omega_nt)$.
Dla całkowitych $r,s\ge0$ pochodna $r$ razy po $x$
i $s$ razy po $t$ wyrazu szeregu jest ograniczona
w module przez
\[
 (n\pi/L)^r\omega_n^s
 \bigl(|a_n|+|b_n|/\omega_n\bigr).
\]
Ponieważ $\omega_n=cn\pi/L$, wybór $m\ge r+s+2$
daje sumowalną majorantę niezależną od $t,x$.
Kryterium Weierstrassa i twierdzenie o różniczkowaniu
szeregów funkcji dają gładkość sumy oraz możliwość
różniczkowania wyraz po wyrazie. Każdy wyraz spełnia
równanie falowe i warunki brzegowe. W chwili $t=0$
szeregi dla $u$ i $u_t$ są odpowiednio szeregami
sinusowymi $g$ i $h$, ponieważ $A_n(0)=a_n$
i $A_n'(0)=b_n$. Z zupełności układu sinusowego
wynika ich zbieżność w $L^2(0,L)$ do $g,h$;
zbieżność jednostajna daje ciągłe sumy, więc równości
z $g,h$ zachodzą w każdym punkcie $[0,L]$.

Dla dowolnego rozwiązania klasy $C^2(\R\times[0,L])$
z zerowymi warunkami Dirichleta różniczkowanie energii
pod znakiem całki i całkowanie przez części dają
\[
 \begin{aligned}
 E'(t)
 &=\operatorname{Re}\int_0^L
 (u_{tt}\overline{u_t}+c^2u_x\overline{u_{xt}})
 \,\mathrm d\lambda_1(x)\\
 &=c^2\operatorname{Re}[u_x\overline{u_t}]_0^L=0.
 \end{aligned}
\]
W drugim kroku użyliśmy równania falowego,
równości pochodnych mieszanych i $u_t(t,0)=u_t(t,L)=0$,
wynikającej z różniczkowania warunków brzegowych po $t$.
Rachunek jest wersją tożsamości energii z rozdziału
~\ref{chap:pde-wave}, poprawną także dla funkcji zespolonych.
Jeżeli $w$ jest różnicą dwóch rozwiązań o tych samych danych,
to $w_t(0,x)=w_x(0,x)=0$, a więc $E_w(0)=0$.
Z zachowania nieujemnej energii i ciągłości pochodnych
wynika $w_t=w_x=0$ dla wszystkich $t,x$. Ponieważ
$w(t,0)=0$, otrzymujemy $w(t,x)=0$.

Możemy również obliczyć energię modalną. Ortogonalność
sinusów oraz kosinusów o dodatnich indeksach daje
\[
 E(t)=\frac L4\sum_{n\ge1}
 \bigl(|A_n'(t)|^2+\omega_n^2|A_n(t)|^2\bigr).
\]
Najpierw stosujemy ortogonalność do sum skończonych,
a potem przechodzimy do granicy dzięki jednostajnej
zbieżności szeregów $u_t,u_x$, a zatem zbieżności w $L^2$.
Rozpisując kwadraty modułów i korzystając z
$\cos^2(\omega_nt)+\sin^2(\omega_nt)=1$, otrzymujemy
$|A_n'|^2+\omega_n^2|A_n|^2
=|b_n|^2+\omega_n^2|a_n|^2$:
wyrazy mieszane z $\operatorname{Re}(a_n\overline{b_n})$
znoszą się. W konsekwencji
\[
 E(t)=\frac L4\sum_{n\ge1}
 (|b_n|^2+\omega_n^2|a_n|^2)=E(0).
\]
Zachowanie energii dotyczy zatem każdego modu osobno,
nie tylko sumy wszystkich modów.
\end{proof}

\begin{przyklad}[Wygładzanie a brak wygładzania]
\label{ex:pde6-smoothing}
Dla równania ciepła amplituda modu $n$ zanika jak
$e^{-\kappa(n\pi/L)^2t}$. Dla równania falowego
współczynnik czasowy modu oscyluje, może chwilowo
znikać, lecz energia nietrywialnego modu pozostaje
dodatnia i stała. Na przykład dla
$g(x)=\sin(n\pi x/L)$, $h=0$ otrzymujemy odpowiednio
$e^{-\kappa\lambda_nt}g(x)$ i
$\cos(c\sqrt{\lambda_n}t)g(x)$.
Chociaż ta funkcja $g$ nie ma zwartego nośnika w $(0,L)$,
oba skończone wzory są rozwiązaniami klasycznymi,
co sprawdzamy przez bezpośrednie różniczkowanie.
Dla równania ciepła w ustalonej chwili $t>0$
mnożnik zanika szybciej niż dowolna potęga $n^{-1}$.
Dla równania falowego w chwili $t=2L/c$ wszystkie
mnożniki $\cos(c\sqrt{\lambda_n}t)$ są równe $1$.
Nie ma więc analogicznego tłumienia wysokich
częstości dla całej ewolucji falowej. Ten kontrast
wyjaśnia spektralnie wygładzanie paraboliczne.
\end{przyklad}

\section{Niejednorodne równania i wzór modalny Duhamela}
\label{sec:pde6-forcing}

Niech $\kappa,L,T>0$. Rozważmy zagadnienie
\[
 \begin{cases}
 u_t-\kappa u_{xx}=f(t,x),&0<t<T,\quad 0<x<L,\\
 u(t,0)=u(t,L)=0,&0\le t\le T,\\
 u(0,x)=g(x),&0\le x\le L.
 \end{cases}
\]
Zakładamy, że $f\in C^\infty([0,T]\times[0,L])$
oraz istnieje zbiór zwarty $K\subset(0,L)$ taki, że
$\operatorname{supp}f(t,\cdot)\subset K$ dla każdego $t\in[0,T]$.
Ponadto niech $g\in C_c^\infty(0,L)$.
Oznaczmy $X_n(x)=\sin(n\pi x/L)$ i $\lambda_n=(n\pi/L)^2$
dla $n\ge1$. Wprowadzamy współczynniki sinusowe
\[
 \begin{aligned}
 f_n(t)&=\frac2L\int_0^L f(t,x)X_n(x)\,\mathrm d\lambda_1(x),\\
 g_n&=\frac2L\int_0^L g(x)X_n(x)\,\mathrm d\lambda_1(x).
 \end{aligned}
\]
Stosujemy więc współczynniki względem funkcji $X_n$,
a nie względem unormowanych funkcji własnych $e_n$.

Najpierw załóżmy, że istnieje rozwiązanie o ciągłych
pochodnych $u_t,u_x,u_{xx}$ na domkniętym prostokącie.
Niech
$b_n(t)=(2/L)\int_0^L u(t,x)X_n(x)\,\mathrm d\lambda_1(x)$.
Ciągłość $u_t$ na zbiorze zwartym pozwala różniczkować
pod znakiem całki. Dwukrotne całkowanie przez części daje
\[
 \begin{aligned}
 \int_0^L u_{xx}X_n\,\mathrm d\lambda_1
 &=[u_xX_n-uX_n']_0^L
   +\int_0^L u X_n''\,\mathrm d\lambda_1\\
 &=-\lambda_n\int_0^L u X_n\,\mathrm d\lambda_1 .
 \end{aligned}
\]
Wyrazy brzegowe znikają, ponieważ zarówno $X_n$,
jak i $u(t,\cdot)$ spełniają zerowe warunki Dirichleta.
Mnożąc równanie przez $(2/L)X_n$ i całkując względem $x$,
otrzymujemy zatem
$b_n'(t)+\kappa\lambda_nb_n(t)=f_n(t)$ i $b_n(0)=g_n$.
Po pomnożeniu przez $e^{\kappa\lambda_nt}$ dostajemy
$(e^{\kappa\lambda_nt}b_n(t))'=e^{\kappa\lambda_nt}f_n(t)$.
Całkowanie od $0$ do $t$ i pomnożenie przez
$e^{-\kappa\lambda_nt}$ prowadzą do wzoru
\begin{equation}\label{eq:pde6-duhamel}
 b_n(t)=e^{-\kappa\lambda_nt}g_n
 +\int_0^t e^{-\kappa\lambda_n(t-s)}f_n(s)
 \,\mathrm d\lambda_1(s).
\end{equation}

Teraz konstruujemy rozwiązanie, nie zakładając jego
istnienia. Definiujemy $b_n$ przez~\eqref{eq:pde6-duhamel}.
Pokażemy, że
$u(t,x)=\sum_{n\ge1}b_n(t)X_n(x)$
jest rozwiązaniem klasy $C^\infty$ na domkniętym
prostokącie (z pochodnymi na brzegu rozumianymi
jako jednostronne lub ich ciągłe przedłużenia).

Dla dowolnych całkowitych $m,j\ge0$ mamy
\[
 |g_n|\le C_m n^{-m},\qquad
 \sup_{t\in[0,T]}|f_n^{(j)}(t)|\le C_{m,j}n^{-m}.
\]
Istotnie, można różniczkować $f_n$ względem czasu
pod znakiem całki. Wszystkie pochodne $f$ i $g$
znikają w ustalonych otoczeniach obu końców odcinka.
Wielokrotne całkowanie przez części względem $x$
nie daje więc wyrazów brzegowych i za każdym razem
wprowadza czynnik $L/(n\pi)$. Normy pochodnych $f$
są ograniczone jednostajnie względem czasu na
domkniętym prostokącie.

Z nieujemności $\kappa\lambda_n$ i $0\le s\le t\le T$
wynika
\[
 \sup_{0\le t\le T}|b_n(t)|
 \le |g_n|+T\sup_{0\le t\le T}|f_n(t)|.
\]
Zatem $b_n$ również zanika szybciej niż dowolna
odwrotna potęga $n$, jednostajnie względem $t$.
Różniczkując~\eqref{eq:pde6-duhamel} według reguły
Leibniza, otrzymujemy $b_n'=f_n-\kappa\lambda_nb_n$.
Dla $j\ge0$ mamy ponadto
$b_n^{(j+1)}=f_n^{(j)}-\kappa\lambda_nb_n^{(j)}$.
Indukcja daje szybkie zanikanie każdej pochodnej
czasowej: jeżeli $b_n^{(j)}$ spełnia oszacowanie
dla każdej potęgi $n^{-m}$, to w wyrazie
$\lambda_nb_n^{(j)}$ stosujemy je z wykładnikiem $m+2$,
ponieważ $\lambda_n=(n\pi/L)^2$.

Dla dowolnych całkowitych $r,j\ge0$ moduł pochodnej
$\partial_x^r\partial_t^j(b_n(t)X_n(x))$ jest
ograniczony przez $(n\pi/L)^r\sup_t|b_n^{(j)}(t)|$.
Wybierając oszacowanie z $m\ge r+2$, otrzymujemy
sumowalną majorantę niezależną od $(t,x)$.
Kryterium Weierstrassa i twierdzenie o różniczkowaniu
szeregów funkcji uzasadniają zbieżność jednostajną
szeregu oraz wszystkich szeregów jego pochodnych.
Zatem
\[
 u_t-\kappa u_{xx}
 =\sum_{n\ge1}(b_n'+\kappa\lambda_nb_n)X_n
 =\sum_{n\ge1}f_n(t)X_n=f(t,\cdot).
\]
Ostatnia równość wynika z zupełności układu sinusowego
(twierdzenie~\ref{thm:pde6-spectrum}) w $L^2(0,L)$
i ciągłości obu stron, gdyż szereg z $f_n(t)$ zbiega
jednostajnie na domkniętym odcinku.
Analogicznie, $b_n(0)=g_n$ i zbieżność szeregu
sinusowego do $g$ w $L^2$ oraz jednostajnie
dają $u(0,x)=g(x)$ dla wszystkich $x\in[0,L]$.
Warunki brzegowe zachodzą wyraz po wyrazie.

Jednoznaczność w klasie rozwiązań ciągłych na domkniętym
prostokącie i klasy $C^{1,2}$ we wnętrzu wynika
z twierdzenia~\ref{thm:pde6-heat}: różnica dwóch rozwiązań
spełnia jednorodne równanie ciepła z zerowymi danymi.
Jeżeli $P_t$ oznacza operator rozwiązania jednorodnego
z sekcji~\ref{sec:pde6-heat}, przy $P_0g=g$, to wzór
modalny można zapisać również jako
\[
 u(t,x)=(P_tg)(x)+
 \int_0^t(P_{t-s}f(s,\cdot))(x)\,\mathrm d\lambda_1(s).
\]
Istotnie, szereg dla $P_{t-s}f(s,\cdot)$ jest jednostajnie
zbieżny dla $0\le s\le t\le T$ i $0\le x\le L$, ponieważ
$e^{-\kappa\lambda_n(t-s)}\le1$ oraz
$\sup_s|f_n(s)|\le C_{2,0}n^{-2}$.
Całkowanie wyraz po wyrazie daje dokładnie
\eqref{eq:pde6-duhamel}. Jest to zasada Duhamela
dla rozważanego zagadnienia Dirichleta.

\section{Regularne zagadnienie Sturma--Liouville'a}
\label{sec:pde6-sturm}

Przy rozdzieleniu zmiennych w równaniach ze
zmiennymi współczynnikami otrzymujemy często operator
$\mathcal L y=-(py')'+qy$. Zdefiniujmy dokładnie
zagadnienie i przestrzeń, w której badamy funkcje własne.

\begin{definicja}[Regularne zagadnienie Dirichleta]
\label{def:pde6-sturm}
Niech $L>0$, $p\in C^1([0,L];\R)$,
$q,w\in C([0,L];\R)$, $p(x)>0$, $w(x)>0$
i $q(x)\ge0$ na $[0,L]$. Zagadnieniem
Sturma--Liouville'a z warunkami Dirichleta nazywamy
\begin{equation}\label{eq:pde6-sturm}
 -(p(x)y'(x))'+q(x)y(x)=\lambda w(x)y(x),
 \quad y(0)=y(L)=0.
\end{equation}
Rozwiązaniem jest $y\in C^2([0,L])$ spełniające
równanie punktowo. Wartością własną nazywamy liczbę
$\lambda\in\mathbb C$, dla której istnieje rozwiązanie
$y\not\equiv0$; funkcje mogą być zespolone. Pracujemy w przestrzeni
$H=L^2((0,L),w(x)\,\mathrm d\lambda_1(x))$ z
iloczynem skalarnym $\langle f,g\rangle_H
=\int_0^Lf(x)\overline{g(x)}w(x)\,\mathrm d\lambda_1(x)$.
Z ciągłości i dodatniości $w$ na zwartym odcinku wynika
$0<\min_{[0,L]}w\le\max_{[0,L]}w<\infty$.
Norma $H$ jest więc równoważna zwykłej normie $L^2$,
a $H$ jest przestrzenią Hilberta.
\end{definicja}

\begin{lemat}[Tożsamość Lagrange'a i rzeczywistość widma]
\label{lem:pde6-lagrange}
Niech $y,z\in C^2([0,L])$. Wtedy
\begin{equation}\label{eq:pde6-lagrange}
 \int_0^L\bigl((\mathcal Ly)(x)\overline{z(x)}
 -y(x)\overline{(\mathcal Lz)(x)}\bigr)\,\mathrm d\lambda_1(x)
 =\left[p\left(y\overline{z'}-y'\overline z\right)\right]_0^L.
\end{equation}
Jeżeli $y\not\equiv0$ jest funkcją własną
z jednorodnymi warunkami Dirichleta, to $\lambda>0$
jest rzeczywista. Funkcje własne dla różnych
wartości własnych są ortogonalne w $H$.
\end{lemat}
\begin{proof}
Całkowanie przez części, przy rzeczywistych $p,q$, daje
\[
 \int_0^L(\mathcal Ly)\overline z\,\mathrm d\lambda_1
 =[-py'\overline z]_0^L
 +\int_0^L(py'\overline{z'}+qy\overline z)\,
 \mathrm d\lambda_1.
\]
Odejmując analogiczną równość dla
$\int_0^L y\overline{\mathcal Lz}\,\mathrm d\lambda_1$,
otrzymujemy~\eqref{eq:pde6-lagrange}.
Dla funkcji własnej $y$ z warunkami Dirichleta
powyższa równość z $z=y$ daje
\[
 \lambda\|y\|_H^2
 =\int_0^L(p(x)|y'(x)|^2+q(x)|y(x)|^2)\,
 \mathrm d\lambda_1(x)>0.
\]
Istotnie, gdyby prawa strona była zerem, dodatniość $p$,
nieujemność $q$ i ciągłość $y'$ wymuszałyby $y'=0$
na całym odcinku, a $y(0)=0$ dawałoby $y=0$.
Ponieważ $\|y\|_H^2>0$, liczba $\lambda$ jest
rzeczywista i dodatnia.
Dla pary funkcji własnych $y,z$ o wartościach
$\lambda,\mu$ stosujemy tożsamość Lagrange'a;
ze znikania brzegów i rzeczywistości liczb własnych
wynika $(\lambda-\mu)\langle y,z\rangle_H=0$.
\end{proof}

\begin{twierdzenie}[Dyskretność i zupełność widma]
\label{thm:pde6-sturm}
Przy założeniach definicji~\ref{def:pde6-sturm}
istnieje rosnący do nieskończoności ciąg
$0<\lambda_1<\lambda_2<\cdots$ i odpowiednie
rzeczywiste funkcje własne $y_n\in C^2([0,L])$,
które po unormowaniu tworzą bazę ortonormalną $H$.
Są to wszystkie wartości własne zagadnienia,
a każda z nich jest prosta.
\end{twierdzenie}
\begin{proof}
Skonstruujemy zwarty operator samosprzężony na $H$,
będący odwrotnością zagadnienia Dirichleta z wagą.
Użyjemy twierdzenia spektralnego dla zwartych
operatorów samosprzężonych w przestrzeni Hilberta:
istnieje baza ortonormalna złożona z wektorów własnych
o niezerowych wartościach własnych oraz bazy jądra;
niezerowe wartości własne mają skończone krotności
i mogą gromadzić się jedynie w zerze.
Jest to jawnie wykorzystywany wynik analizy funkcjonalnej.

\emph{Konstrukcja funkcji Greena.}
Niech $v_1,v_2$ będą rzeczywistymi rozwiązaniami
$\mathcal Lv=0$ o danych
$v_1(0)=0$, $v_1'(0)=1$ oraz
$v_2(L)=0$, $v_2'(L)=-1$.
Istnienie i jednoznaczność wynikają z twierdzenia
dla liniowych układów ODE o ciągłych współczynnikach:
równanie zapisujemy jako układ $v'=z/p$, $z'=qv$,
gdzie $z=pv'$.
Nietrywialne rozwiązanie równania jednorodnego
nie może znikać na obu końcach, ponieważ rachunek
energii z lematu~\ref{lem:pde6-lagrange} przy $\lambda=0$
wymusza wtedy $v=0$.
W szczególności $v_1(L)\ne0$.
Z równań $(pv_1')'=qv_1$, $(pv_2')'=qv_2$ wynika
\[
 [p(v_1'v_2-v_1v_2')]'=qv_1v_2-v_1qv_2=0.
\]
Zatem $C=p(s)(v_1'(s)v_2(s)-v_1(s)v_2'(s))$
jest stałą, a obliczenie w $s=L$ daje
$C=p(L)v_1(L)\ne0$.
Definiujemy
\begin{equation}\label{eq:pde6-green}
 G(x,s)=\frac{v_1(\min\{x,s\})v_2(\max\{x,s\})}{C}.
\end{equation}
Jądro jest rzeczywiste, ciągłe i symetryczne na
$[0,L]^2$ oraz $G(0,s)=G(L,s)=0$.
Dla $s\in(0,L)$ spełnia równanie jednorodne po obu
stronach punktu $x=s$. W tym punkcie zachodzi
\[
 -p(s)(\partial_xG(s^+,s)-\partial_xG(s^-,s))
 =\frac{p(s)(v_1'(s)v_2(s)-v_1(s)v_2'(s))}{C}=1.
\]
Jest to warunek skoku funkcji Greena, wyrażający
$\mathcal L_xG(x,s)=\delta_s$ we wnętrzu odcinka.
Nie potrzebujemy tu ogólnej teorii dystrybucji:
własność odwrotności sprawdzimy bezpośrednio.

\emph{Ograniczoność, zwartość i samosprzężoność.}
Niech $\mathrm d\nu(s)=w(s)\,\mathrm d\lambda_1(s)$
oraz $M=\nu((0,L))>0$. Definiujemy
$Tf(x)=\int_0^LG(x,s)f(s)\,\mathrm d\nu(s)$.
Ponieważ $\|f\|_{L^1(\nu)}\le\sqrt M\|f\|_H$,
całka jest bezwzględnie zbieżna dla każdego $x$,
nie zależy od reprezentanta $f$, a jednostajna
ciągłość $G$ daje ciągłość $Tf$ względem $x$.
Ponadto $\|Tf\|_\infty\le\|G\|_\infty\sqrt M\|f\|_H$,
skąd $\|Tf\|_H\le M\|G\|_\infty\|f\|_H$.

Wybieramy podziały $x_j=jL/N$, $j=0,\ldots,N$,
oraz ciągłe odcinkami liniowe funkcje kapeluszowe
$\theta_j$ przyjmujące wartości $\delta_{ij}$ w węzłach $x_i$.
Są nieujemne, sumują się do $1$, a $\theta_j(x)>0$
implikuje $|x-x_j|\le L/N$.
Interpolant
$G_N(x,s)=\sum_{i,j=0}^NG(x_i,x_j)\theta_i(x)\theta_j(s)$
jest kombinacją wypukłą wartości $G$ w węzłach
odległych od $(x,s)$ o co najwyżej $\sqrt2 L/N$.
Z jednostajnej ciągłości $G$ wynika
$\|G_N-G\|_\infty\to0$.
Operator $T_N$ o jądrze $G_N$ ma obraz zawarty
w przestrzeni rozpiętej przez $\theta_0,\ldots,\theta_N$,
a więc ma skończony rząd. Powyższe oszacowanie
zastosowane do różnicy jąder daje
$\|T_N-T\|_{H\to H}\le M\|G_N-G\|_\infty\to0$.
Korzystamy z faktu, że granica operatorów zwartych
w normie operatorowej jest zwarta, więc $T$ jest zwarty.

Dla $f,h\in H$ całka podwójna z modułu
$G(x,s)f(s)\overline{h(x)}$ względem $\nu\otimes\nu$
jest skończona, bo nie przekracza
$\|G\|_\infty\|f\|_{L^1(\nu)}\|h\|_{L^1(\nu)}$.
Z twierdzenia Fubiniego i rzeczywistej symetrii jądra
otrzymujemy
\[
 \langle Tf,h\rangle_H
 =\int_0^L\int_0^L G(x,s)f(s)\overline{h(x)}
 \,\mathrm d\nu(s)\,\mathrm d\nu(x)
 =\langle f,Th\rangle_H.
\]
Operator jest ograniczony i określony na całym $H$,
a więc jest samosprzężony.

\emph{Regularność i własność odwrotności.}
Dla $f\in H$ funkcja $F=wf$ należy do $L^1(0,L)$.
Wprowadzamy absolutnie ciągłe funkcje
\[
 A(x)=\int_0^x v_1(s)F(s)\,\mathrm d\lambda_1(s),
 \qquad
 B(x)=\int_x^L v_2(s)F(s)\,\mathrm d\lambda_1(s).
\]
Mamy $A'=v_1F$, $B'=-v_2F$ prawie wszędzie.
Z definicji jądra funkcja $y=Tf$ ma postać
$y=(v_2 A+v_1 B)/C$, zatem jest absolutnie ciągła
i znika w $0,L$.
Różniczkowanie prawie wszędzie daje
$y'=(v_2'A+v_1'B)/C$, ponieważ składniki z $A',B'$
znoszą się. Prawa strona jest ciągła, więc z twierdzenia
podstawowego rachunku całkowego dla funkcji absolutnie
ciągłych wynika, że $y$ jest klasy $C^1$ i wzór na $y'$
zachodzi wszędzie (na końcach jednostronnie).
Funkcja $py'=(pv_2'A+pv_1'B)/C$ jest absolutnie ciągła.
Różniczkując ją prawie wszędzie i używając równań
jednorodnych dla $v_1,v_2$, dostajemy
\[
 \begin{aligned}
 (py')'
 &=q(v_2A+v_1B)/C\\
 &\quad +p(v_2'v_1-v_1'v_2)F/C
 =qy-F.
 \end{aligned}
\]
Zatem $\mathcal L(Tf)=wf$ prawie wszędzie.
Jeśli $F$ ma ciągły reprezentant, prawa strona wzoru
na $(py')'$ jest ciągła, więc $py'\in C^1$, a
$p\in C^1$ i $p>0$ dają $y\in C^2$.
Dla dowolnego $f\in H$ nie twierdzimy jeszcze,
że $Tf$ jest klasy $C^2$.
Dla ciągłego wymuszenia $F$ powyższa konstrukcja
z $f=F/w$ daje rozwiązanie klasy $C^2$ równania
$\mathcal Ly=F$; jednoznaczność wynika z rachunku energii
dla równania jednorodnego.

\emph{Dodatniość i przejście do wartości własnych.}
Całkowanie przez części jest dozwolone, ponieważ
$py'$ jest absolutnie ciągła, $y\in C^1$ i $y(0)=y(L)=0$.
Z równania dla $y=Tf$ wynika
\[
 \langle Tf,f\rangle_H
 =\int_0^L y\overline{\mathcal Ly}\,\mathrm d\lambda_1
 =\int_0^L(p|y'|^2+q|y|^2)\,\mathrm d\lambda_1\ge0.
\]
Jeśli energia jest zerowa, ciągłość $y'$ i dodatniość $p$
dają $y'=0$, a warunki brzegowe dają $y=0$.
Wtedy $\mathcal Ly=wf=0$ prawie wszędzie, czyli $f=0$.
Zatem $T$ jest ściśle dodatni i ma trywialne jądro.

Twierdzenie spektralne daje bazę ortonormalną
wektorów własnych $Ty_n=\mu_n y_n$, gdzie $\mu_n>0$.
Możemy wybrać rzeczywiste wektory własne, ponieważ
rzeczywistość jądra daje $T\overline f=\overline{Tf}$.
Istotnie, rzeczywista i urojona część każdego wektora
własnego należą do tej samej przestrzeni własnej;
ortonormalizację można przeprowadzić na jej części
rzeczywistej, która rozpina ją nad $\mathbb C$.
Każda przestrzeń własna odpowiadająca $\mu_n\ne0$
jest skończenie wymiarowa.
Przestrzeń $H$ jest nieskończenie wymiarowa (na przykład
zawiera dowolnie wiele liniowo niezależnych funkcji
charakterystycznych rozłącznych podprzedziałów).
Ponieważ jądro jest trywialne, niezerowych wartości
własnych jest nieskończenie wiele i można je uporządkować
nierosnąco z krotnościami, $\mu_n\downarrow0$.

Wybierając reprezentanta $y_n=\mu_n^{-1}Ty_n$,
otrzymujemy $y_n\in C^1([0,L])$ i warunki Dirichleta.
Równanie $\mathcal L(Ty_n)=wy_n$ daje
$\mathcal Ly_n=\mu_n^{-1}wy_n$ prawie wszędzie.
Wymuszenie $\mu_n^{-1}wy_n$ jest ciągłe, więc
z wykazanej regularności wynika $y_n\in C^2([0,L])$
oraz równanie punktowe. Stąd $\lambda_n=\mu_n^{-1}$
są wartościami własnymi rozważanego zagadnienia.
Odwrotnie, jeśli $\mathcal Ly=\lambda wy$ dla
nietrywialnego $y$ z warunkami Dirichleta, to
$\lambda>0$ z lematu~\ref{lem:pde6-lagrange}.
Funkcje $T(\lambda y)$ i $y$ rozwiązują to samo
zagadnienie z ciągłym wymuszeniem $\lambda wy$;
jego jednoznaczność daje $Ty=\lambda^{-1}y$.
Nie ma więc innych wartości własnych.

Dla ustalonej $\lambda$ rozwiązanie równania
$(py')'=(q-\lambda w)y$ z $y(0)=0$ jest jednoznacznie
wyznaczone przez $y'(0)$, z twierdzenia o liniowych
układach ODE. Gdyby również $y'(0)=0$, rozwiązanie
byłoby zerowe. Dwa nietrywialne rozwiązania z zerowym
warunkiem w $0$ są więc proporcjonalne, co daje
prostotę wartości własnych. Ostatecznie możemy
uporządkować je ściśle rosnąco,
$0<\lambda_1<\lambda_2<\cdots\to\infty$,
a odpowiadające im unormowane funkcje własne
tworzą bazę ortonormalną $H$.
\end{proof}

\begin{uwaga}[Co jest tu nowe?]
Dla stałych współczynników mieliśmy jawne funkcje
$\sin(n\pi x/L)$; przy zmiennych współczynnikach
możemy zwykle nie znać ich w postaci elementarnej,
ale nadal mamy dyskretne wartości własne oraz
zupełny układ modów. To właśnie uzasadnia metodę
rozdzielenia zmiennych również bez jawnych wzorów.
Twierdzenie o oscylacji Sturma, mówiące o liczbie
zer kolejnych funkcji własnych, omówimy osobno,
gdy będziemy dysponować szerszą teorią ODE i
zagadnień spektralnych.
\end{uwaga}

\section{Wyższe wymiary: prostokąt i geometria widma}
\label{sec:pde6-rectangle}

Niech $L_1,L_2>0$ i $\Omega=(0,L_1)\times(0,L_2)$.
Rozważamy klasyczne zagadnienie własne Dirichleta
$-\Delta X=\lambda X$ w $\Omega$, $X|_{\partial\Omega}=0$,
dla $X\in C^2(\overline\Omega)$.
Pojedyncze rozwiązania rozdzielone szukamy w postaci
$X(x,y)=U(x)V(y)$. Jeśli iloczyn jest nietrywialny,
to wybór punktu $y_*$, w którym $V(y_*)\ne0$, w równaniu
$-U''V-UV''=\lambda UV$ daje
$-U''=\alpha U$, gdzie
$\alpha=\lambda+V''(y_*)/V(y_*)$.
Następnie wybór $x_*$ z $U(x_*)\ne0$ daje
$-V''=\beta V$, gdzie $\beta=\lambda-\alpha$.
Warunki brzegowe wymuszają
$U(0)=U(L_1)=V(0)=V(L_2)=0$.
Z twierdzenia~\ref{thm:pde6-spectrum} wynika więc,
że z dokładnością do stałego czynnika otrzymujemy
\begin{equation}\label{eq:pde6-rectangle}
 X_{mn}(x,y)=\sin\left(\frac{m\pi x}{L_1}\right)
 \sin\left(\frac{n\pi y}{L_2}\right),\quad
 \lambda_{mn}=\pi^2\left(\frac{m^2}{L_1^2}
 +\frac{n^2}{L_2^2}\right),\quad m,n\ge1.
\end{equation}
Odwrotnie, każda z tych funkcji spełnia równanie
i wszystkie warunki brzegowe, co wynika bezpośrednio
z różniczkowania obu czynników. Pozostaje wykazać,
że rozdzielenie zmiennych nie pomija żadnych wartości
własnych ani funkcji własnych.

Niech $e_m^{(1)}$ i $e_n^{(2)}$ oznaczają unormowane
funkcje sinusowe na odcinkach $(0,L_1)$ i $(0,L_2)$:
\[
 e_m^{(1)}(x)=\sqrt{2/L_1}\sin(m\pi x/L_1),\qquad
 e_n^{(2)}(y)=\sqrt{2/L_2}\sin(n\pi y/L_2).
\]
Piszemy $E_{mn}(x,y)=e_m^{(1)}(x)e_n^{(2)}(y)$,
czyli $E_{mn}=2(L_1L_2)^{-1/2}X_{mn}$.
Z twierdzenia Fubiniego i ortonormalności baz
jednowymiarowych otrzymujemy
\[
 \langle E_{mn},E_{ij}\rangle_{L^2(\Omega)}
 =\delta_{mi}\delta_{nj}.
\]
Tutaj całkujemy względem zwykłej miary Lebesgue'a
$\lambda_2$ na prostokącie.

Wykażemy zupełność tej rodziny. Niech $f\in L^2(\Omega)$
będzie ortogonalna do wszystkich $E_{mn}$.
Z twierdzenia Fubiniego dla prawie każdego $y$
funkcja $f(\cdot,y)$ należy do $L^2(0,L_1)$.
Dla $m\ge1$ określamy prawie wszędzie
\[
 F_m(y)=\int_0^{L_1} f(x,y)\overline{e_m^{(1)}(x)}
 \,\mathrm d\lambda_1(x).
\]
Z nierówności Cauchy'ego--Schwarza i unormowania
$e_m^{(1)}$ wynika
$|F_m(y)|^2\le\int_0^{L_1}|f(x,y)|^2\,\mathrm d\lambda_1(x)$.
W szczególności $F_m\in L^2(0,L_2)$.
Iloczyn $f\overline{E_{mn}}$ jest całkowalny na
prostokącie, ponownie z nierówności Cauchy'ego--Schwarza.
Z twierdzenia Fubiniego i założonej ortogonalności mamy
\[
 \langle F_m,e_n^{(2)}\rangle_{L^2(0,L_2)}
 =\langle f,E_{mn}\rangle_{L^2(\Omega)}=0
 \quad(n\ge1).
\]
Z zupełności bazy sinusowej na $(0,L_2)$ wynika
$F_m=0$ prawie wszędzie. Ponieważ zbiór indeksów $m$
jest przeliczalny, istnieje jeden zbiór pełnej miary
w $(0,L_2)$, na którym wszystkie $F_m(y)$ znikają
i $f(\cdot,y)\in L^2(0,L_1)$.
Dla każdego takiego $y$ zupełność bazy na $(0,L_1)$
daje $f(\cdot,y)=0$ w $L^2(0,L_1)$.
Ponowne zastosowanie twierdzenia Fubiniego daje
$\|f\|_{L^2(\Omega)}=0$. Dopełnienie ortogonalne
rozpatrywanej rodziny jest więc trywialne, co
dowodzi, że $(E_{mn})_{m,n\ge1}$ jest bazą ortonormalną.

Niech teraz $X\not\equiv0$ będzie dowolną klasyczną
funkcją własną z warunkami Dirichleta.
Dwukrotne całkowanie przez części osobno względem
każdej ze zmiennych, przy ustalonej drugiej zmiennej,
daje
\[
 \langle-\Delta X,E_{mn}\rangle
 =\langle X,-\Delta E_{mn}\rangle.
\]
Wyrazy brzegowe znikają, ponieważ $X$ oraz $E_{mn}$
zerują się na wszystkich bokach prostokąta.
Stąd $(\lambda-\lambda_{mn})\langle X,E_{mn}\rangle=0$
dla wszystkich $m,n$. Co najmniej jeden współczynnik
jest niezerowy z zupełności bazy, więc $\lambda$
jest jedną z liczb~\eqref{eq:pde6-rectangle}.
Dla ustalonego $\lambda$ istnieje tylko skończenie
wiele par $(m,n)$ z $\lambda_{mn}=\lambda$, gdyż
$m\le L_1\sqrt\lambda/\pi$ i $n\le L_2\sqrt\lambda/\pi$.
Rozwinięcie $X$ w bazie zawiera wyłącznie te wyrazy,
a równość w $L^2$ jest równością punktową dzięki
ciągłości. Zatem przestrzeń własna jest dokładnie
rozpięta przez $E_{mn}$ odpowiadające tym parom,
a jej wymiar jest liczbą takich par.
W szczególności w każdym ograniczonym przedziale
liczbowym jest tylko skończenie wiele wartości własnych.

\begin{uwaga}[Geometria prostokąta a krotności]
Dla kwadratu $L_1=L_2=L$ mamy na przykład
$\lambda_{12}=\lambda_{21}=5\pi^2/L^2$.
Są to jedyne dodatnie pary całkowite z $m^2+n^2=5$,
więc ta wartość własna ma krotność dokładnie $2$.
Jeżeli natomiast $L_1^2/L_2^2$ jest niewymierne,
to z równości $\lambda_{mn}=\lambda_{ij}$ wynika
\[
 m^2-i^2=\frac{L_1^2}{L_2^2}(j^2-n^2).
\]
Lewa strona jest całkowita. Gdyby $j^2-n^2\ne0$,
stosunek $L_1^2/L_2^2$ byłby wymierny. Zatem
$j=n$, a następnie $m=i$, bo wszystkie indeksy
są dodatnie. W tym przypadku każda wartość własna
jest prosta.
\end{uwaga}

\section{Dlaczego nie można bezwarunkowo odwrócić ewolucji cieplnej?}
\label{sec:pde6-backward}

Niech $\kappa,L,T>0$ i $H=L^2(0,L)$ z normą względem
miary Lebesgue'a. Rozważmy operator $P_T:H\to H$,
przyporządkowujący danym początkowym $g$ rozwiązanie
równania ciepła z jednorodnymi warunkami Dirichleta
w chwili $T$. Z twierdzenia~\ref{thm:pde6-spectrum}
i stwierdzenia~\ref{prop:pde6-L2-semigroup} wynika
\[
 P_Tg=\sum_{n\ge1} e^{-\kappa\lambda_nT}g_ne_n,\qquad
 g_n=\langle g,e_n\rangle_H,\qquad
 \lambda_n=(n\pi/L)^2.
\]
Szereg zbiega w $H$, a operator jest liniowy i ograniczony:
z tożsamości Parsevala wynika
$\|P_Tg\|_H\le e^{-\kappa\lambda_1T}\|g\|_H$.

Rozważmy zagadnienie odwrotne: dla danych końcowych
$h\in H$ znaleźć $g\in H$ takie, że $P_Tg=h$.
Należy rozróżnić trzy aspekty poprawności zagadnienia
w sensie Hadamarda.

\emph{Jednoznaczność.}
Wszystkie mnożniki $e^{-\kappa\lambda_nT}$ są dodatnie.
Jeżeli $P_Tg=0$, to każdy współczynnik $g_n$ jest zerowy,
a zupełność bazy daje $g=0$. Operator $P_T$ jest zatem
iniektywny, więc dla każdego $h$ z jego obrazu dane
początkowe są wyznaczone jednoznacznie.

\emph{Istnienie.}
Dla $h_n=\langle h,e_n\rangle_H$ warunek $P_Tg=h$
wymusza $g_n=e^{\kappa\lambda_nT}h_n$.
Z Parsevala wynika warunek konieczny
\[
 \sum_{n\ge1} e^{2\kappa\lambda_nT}|h_n|^2<\infty.
\]
Jest on również wystarczający: jeśli zachodzi,
to ortonormalność i zupełność przestrzeni Hilberta
zapewniają zbieżność szeregu
$g=\sum_{n\ge1}e^{\kappa\lambda_nT}h_ne_n$ w $H$.
Wstawienie jego współczynników do wzoru na $P_Tg$
daje $P_Tg=h$. Tym samym
\[
 \operatorname{ran}P_T
 =\left\{h\in H:
 \sum_{n\ge1}e^{2\kappa\lambda_nT}|h_n|^2<\infty\right\}.
\]
Obraz jest gęsty w $H$, ponieważ zawiera każdą skończoną
kombinację liniową funkcji $e_n$.
Nie jest jednak całym $H$. Istotnie, szereg
$h_*=\sum_{n\ge1}e^{-\kappa\lambda_nT}e_n$
zbiega w $H$, ponieważ $\lambda_n$ rośnie jak $n^2$,
ale dla jego współczynników powyższa suma
ważona jest równa $\sum_{n\ge1}1=\infty$.
Zatem $h_*\notin\operatorname{ran}P_T$.
W szczególności obraz $P_T$ nie jest domknięty.

\emph{Stabilność.}
Nawet jeśli ograniczymy dane końcowe do
$\operatorname{ran}P_T$, odwrotność nie jest ciągła
w normie $H$ odziedziczonej z przestrzeni danych.
Niech $h^{(n)}=P_Te_n=e^{-\kappa\lambda_nT}e_n$.
Wówczas
\[
 \|h^{(n)}\|_H=e^{-\kappa\lambda_nT}\longrightarrow0,
 \qquad
 \|P_T^{-1}h^{(n)}\|_H=\|e_n\|_H=1.
\]
Ponieważ $P_T^{-1}0=0$, przeczy to ciągłości
odwrotności w zerze. Przez liniowość odwrotność nie
jest ciągła w żadnym punkcie swojej dziedziny $\operatorname{ran}P_T$.
Dowolnie małe błędy danych końcowych mogą zatem
odpowiadać błędom danych początkowych, które nie
dążą do zera. Jest to spektralna wersja
niestabilności Hadamarda z rozdziału~\ref{chap:pde-heat}.

\section{Podsumowanie i dalsza perspektywa}
\label{sec:pde6-summary}

Dzięki średnim Fejéra uzyskaliśmy zupełność układu
trygonometrycznego i twierdzenie Parsevala.
Na ich podstawie skonstruowaliśmy jawne rozwiązania
ciepła i fal na odcinku oraz ostre oszacowanie zaniku.
Metoda ta przechodzi na zmienne współczynniki przez
regularną teorię Sturma--Liouville'a i na prostokąty
przez iloczyny tensorowe. W kolejnych rozdziałach
rozszerzymy ją na rozwiązania słabe, operatory
samosprzężone, metody wariacyjne i transformatę
Fouriera na $\R^n$.


